\documentclass[12pt]{article}

\input{../../../_assets/preambulo.tex}


\begin{document}

    % 1. Foto de fondo
    % 2. Título
    % 3. Encabezado Izquierdo
    % 4. Color de fondo
    % 5. Coord x del titulo
    % 6. Coord y del titulo
    % 7. Fecha

    
    \input{../../../_assets/portada}
    \portadaExamenFotoDif{../../dueA4.jpg}{Internet of\\Things}{Internet of Things}{MidnightBlue}{-8}{28}{2026}{Arturo Olivares Martos}

    \begin{description}
        \item[Asignatura] Internet of Things.
        \item[Curso Académico] Winter Semester 2021.
        % \item[Grado] ---.
        % \item[Grupo] ---.
        \item[Profesor] Gregor Schiele.
        \item[Descripción] Klausur WS.
        \item[Fecha] 04/03/2021.
        \item[Duración] 90 minutes.
    
    \end{description}
    \newpage


    % ------------------------------------
    
    \begin{ejercicio}[20 points]
        ~
        \begin{enumerate}[label=\alph*)]
            \item An \verb|802.15.4| network is deployed in the following three application scenarios. Which CSMA/CA MAC Layer version would you use? Justify your answer.
            \begin{enumerate}[label=\roman*)]
                \item (3 points) A smart home network, where only $10$ temperature sensors are installed. Each sensor reports a measurement once per minute.
                \item (3 points) A smart clinic network that transmits monitoring data about the patients and alarm notifications.
                \item (3 points) A smart farm network, where the sensor nodes are distributed over a large area. Some coordinators are powered by solar battery.
            \end{enumerate}
            \item (4 points) When unslotted CSMA/CA is applied, at each trail the back-off window size will be doubled. However, it doesn't always lead to a longer delay, why?
            \item (3 points) Name one of the reasons why we need 6LoWPAN to support IPv6 over \verb|802.15.4|?
            \item (4 points) A single IPv6 packet with a size of $555$ Byte (including the IPv6 header and its payload) is sent to an IoT device using \verb|802.15.4| and 6LoWPAN. The IPv6 header is uncompressed and we use 6LoWPAN fragmentation. What would be the value of the datagram offset field in the 5th fragment? Justify your answer.
            
            \begin{observacion}
            We assume each fragment has $72$ bytes available for its payload.
            \end{observacion}
        \end{enumerate}
    \end{ejercicio}

    \newpage

    \begin{ejercicio}[15 points]
        ~
        \begin{enumerate}[label=\alph*)]
            \item Given the following three application scenarios, should we apply NB-IoT or \verb|802.15.4| as the networking technique? Justify your answer.
            \begin{enumerate}[label=\roman*)]
                \item (3 points) We want to deploy a people counting sensor at the front door of each university building. The sensor data should be sent to the data sink periodically.
                \item (3 points) We want to develop a sheep fence with motion sensors for a farm of about $900$ square metres. We want to get an alarm/notification when a sheep tries to escape.
                \item (3 points) A bicycle sharing company in Duisburg wants to develop a bicycle tracking system. Bicycles should be tracked at all times, including when they are left in an underground parking lot.
            \end{enumerate}
            \item (6 points) Provide the names of the two different types of channels used in NB-IoT, and explain what each of them are used for. On which type of channel do we send Paging messages?
        \end{enumerate}
    \end{ejercicio}

    \newpage

    \begin{ejercicio}[20 points]
        In the lecture you learned about the application layer protocol MQTT.
        \begin{enumerate}[label=\alph*)]
            \item (4 points) Describe a scenario in which you would apply the retain flag of MQTT. Justify your answer.
            \item Given an MQTT configuration with four clients C1, C2, C3, C4 and a broker B:
            \begin{enumerate}[label=\roman*)]
                \item (3 points)
                \begin{itemize}
                    \item C1 publishes a message to topic \verb|basement/temperature|,
                    \item C2 publishes a message to topic \verb|groundfloor/kitchen/temperature|,
                    \item C3 publishes a message to topic \verb|secondfloor/bathroom/temperature|.
                \end{itemize}
                C4 should subscribe to which topic to receive all messages from C1, C2 and C3?
                \item (3 points)
                \begin{itemize}
                    \item C1 publishes a message to topic \verb|bathroom/ceiling/temp|,
                    \item C2 publishes a message to topic \verb|bathroom/floor/temp|,
                    \item C3 publishes a message to topic \verb|bathroom/boiler/outlet/temp|.
                \end{itemize}
                C4 should subscribe to which topic to get the environment temperature of the bathroom?
                \begin{observacion}
                    We do not want to receive any other types of measurements, e.g. humidity.
                \end{observacion}
                \item (3 points) C1, C2, C3 subscribe to shared subscriptions \verb|\$share/exam/iot/temp|. To which topic should C4 publish, such that C1, C2, or C3 will receive the message?
            \end{enumerate}
            \item (3 points) MQTT v5 introduced multiple new features, among them ``Topic Alias''. While creating a ``Topic Alias'', why should the central broker store the topic alias mapping?
            
            \item (4 points) Client C2 subscribes to the topic \verb|a/b/c| with QoS level $0$. Afterwards, client C1 publishes a message to the topic \verb|a/b/c| with QoS level $2$. Looking at the message headers: what is the QoS level between C1 and the broker? What is the QoS level between the broker and C2?
        \end{enumerate}
    \end{ejercicio}

    \newpage

    \begin{ejercicio}[14 points]
        ~
        \begin{enumerate}[label=\alph*)]
            \item\label{itm:robot1}
                (6 points) A robot moves linearly with speed $v = \unitfrac[2]{m}{s}$. It starts at time $t = 0$ at location $m_0 = \unit[0]{m}$. We want it to send us its location with the ``Distance-based Reporting'' protocol from the lecture with an accuracy of $\text{acc} = \unit[10]{m}$ and max speed $v_{\max} = \unitfrac[4]{m}{s}$. Write down the next three positions after $m_0$ ($m_1, m_2, m_3$) that the robot will report.
            \item (4 points) When the robot from Task~\ref{itm:robot1} keeps moving in the same direction with the same speed for a longer time, ``Distance-based Reporting'' leads to energy inefficiencies. Can we apply ``Dead Reckoning'' to make it more efficient? Justify your answer.
            \item (4 points) Describe a scenario where ``Dead Reckoning'' leads to energy inefficiency. Justify your answer.
        \end{enumerate}
    \end{ejercicio}

    \newpage

    \begin{ejercicio}[11 points]
        ~
        \begin{enumerate}[label=\alph*)]
            \item In the following two IoT application scenarios artificial intelligence shall be used to solve a task. Would you deploy the AI (or a specific part of it) locally on an embedded device or in the cloud? Justify your answer.
            \begin{enumerate}[label=\roman*)]
                \item (3 points) A voice-controlled virtual assistant (such as Amazon Alexa) shall wake up from a low-energy state on recognizing the user saying a specific keyword.
                \item (3 points) Imagine a video surveillance system for intruder detection that includes a smart camera. If a suspicious object is detected in the scene, it must be analysed whether it actually is an intruder.
            \end{enumerate}
            \item (2 points) Explain the difference between pruning and quantization of an artificial neural network.
            \item (3 points) How does magnitude-based pruning decide which weights to prune? Why can that decision lead to drastic performance drops?
        \end{enumerate}
    \end{ejercicio}

    \newpage
    \setcounter{ejercicio}{0} % Reiniciar contador de ejercicios
    % \noindent
    % \textbf{Solución.}
    \newpage













    \begin{ejercicio}[20 points]
        ~
        \begin{enumerate}[label=\alph*)]
            \item An \verb|802.15.4| network is deployed in the following three application scenarios. Which CSMA/CA MAC Layer version would you use? Justify your answer.
            \begin{enumerate}[label=\roman*)]
                \item (3 points) A smart home network, where only $10$ temperature sensors are installed. Each sensor reports a measurement once per minute.
                
                Since only $10$ sensors are installed and they report their measurements once per minute, the network is not very dense and the traffic is really low. Therefore, a network without Beacons, which uses unslotted CSMA/CA, is sufficient for this application scenario.
                
                Having to wait for a Congestion Window to be opened (in the case of slotted CSMA/CA) and adding the complexity of the Beacons is not necessary in this case, since the network is not dense and the traffic is low. In addition, since they are installed in a home, those sensors may be plugged to the power supply, so they do not have energy constraints.
                
                
                \item (3 points) A smart clinic network that transmits monitoring data about the patients and alarm notifications.
                
                In this case, the network is dense and the traffic is high. In addition, given the importance of the alarm notifications, they need to be garanteed to be able to use the chanel. Therefore, a network with Beacons, which uses slotted CSMA/CA, is the best option for this application scenario.

                Those alarm notifications could be sent using the Guaranteed Time Slots (GTS) of the Contention Free Period (CFP) of the Beacons, which allows to guarantee the access to the channel for those alarm notifications.
                \item (3 points) A smart farm network, where the sensor nodes are distributed over a large area. Some coordinators are powered by solar battery.
                
                In this case, the main aspect is the energy consumption, since some coordinators are powered by solar battery. Therefore, a network with Beacons would be the best option, where the PAN Coordinator could increase the Inactive Period of the frames, so that the coordinators can sleep for longer periods of time and save energy.
            \end{enumerate}
            \item (4 points) When unslotted CSMA/CA is applied, at each trail the back-off window size will be doubled. However, it doesn't always lead to a longer delay, why?
            
            In each trial, the back-off period is randomly selected from the back-off window. Even though the back-off window size is doubled, it is possible that the randomly selected back-off period is smaller than the previous one, leading to a shorter delay.
            \item (3 points) Name one of the reasons why we need 6LoWPAN to support IPv6 over \verb|802.15.4|?
            
            According to the IPv6 specification, the minimum MTU size for IPv6 is $1280$ bytes. If the MTU size of the link layer is smaller than $1280$ bytes, the IPv6 specification requires fragmentation in the link layer. Since the MTU size of \verb|802.15.4| is $127$ bytes, we need 6LoWPAN to handle that fragmentation and reassembly of IPv6 packets over \verb|802.15.4|.
            \item (4 points) A single IPv6 packet with a size of $555$ Byte (including the IPv6 header and its payload) is sent to an IoT device using \verb|802.15.4| and 6LoWPAN. The IPv6 header is uncompressed and we use 6LoWPAN fragmentation. What would be the value of the datagram offset field in the 5th fragment? Justify your answer.
            
            \begin{observacion}
            We assume each fragment has $72$ bytes available for its payload.
            \end{observacion}

            Let's calculate the total offset of each fragment. Since $72$ is divisible by $8$, all the available payload bytes of each fragment will be used.
            \begin{itemize}
                \item The first fragment has an offset of $0$ bytes. It carries the bytes $0-71$ of the original packet.
                \item The second fragment has an offset of $72$ bytes. It carries the bytes $72-143$ of the original packet.
                \item The third fragment has an offset of $144$ bytes. It carries the bytes $144-215$ of the original packet.
                \item The fourth fragment has an offset of $216$ bytes. It carries the bytes $216-287$ of the original packet.
                \item The fifth fragment has an offset of $288$ bytes. It carries the bytes $288-359$ of the original packet.
            \end{itemize}

            Therefore, the offset field in the 5th fragment would be $\nicefrac{288}{8} = 36$.
        \end{enumerate}
    \end{ejercicio}

    \newpage

    \begin{ejercicio}[15 points]
        ~
        \begin{enumerate}[label=\alph*)]
            \item Given the following three application scenarios, should we apply NB-IoT or \verb|802.15.4| as the networking technique? Justify your answer.
            \begin{enumerate}[label=\roman*)]
                \item (3 points) We want to deploy a people counting sensor at the front door of each university building. The sensor data should be sent to the data sink periodically.
                
                Since the university buildings are located in an small area, the sensor nodes can be connected to a single \verb|802.15.4| network, which is a low-power and low-cost solution. Since it is in the license-free spectrum, it is much cheaper than using NB-IoT, which requires a subscription to a mobile network operator. Therefore, \verb|802.15.4| is the best option for this application scenario.

                A possible drawback could be the distance between the buildings, since \verb|802.15.4| has a limited range. However, using a peer-to-peer topology, each sensor will be able to reach the data sink, even if it is located in a different building.
                \item (3 points) We want to develop a sheep fence with motion sensors for a farm of about $900$ square metres. We want to get an alarm/notification when a sheep tries to escape.
                
                Since the farm is located in a small area, the sensor nodes can be connected to a single \verb|802.15.4| network, which is a low-power and low-cost solution. Since it is in the license-free spectrum, it is much cheaper than using NB-IoT, which requires a subscription to a mobile network operator. Therefore, \verb|802.15.4| is the best option for this application scenario.
                \item (3 points) A bicycle sharing company in Duisburg wants to develop a bicycle tracking system. Bicycles should be tracked at all times, including when they are left in an underground parking lot.
                
                Since the bicycles can be located anywhere in Duisburg, \verb|802.15.4| is not a suitable option for this application scenario, since it has a limited range and could not reach all the bicycles. Therefore, NB-IoT is the best option for this application scenario, since it uses the mobile network infrastructure.

                In addition, since the bicycles can be located in an underground parking lot, NB-IoT is a better option than \verb|802.15.4|, since it has better penetration in buildings and underground areas.
            \end{enumerate}
            \item (6 points) Provide the names of the two different types of channels used in NB-IoT, and explain what each of them are used for. On which type of channel do we send Paging messages?
            
            There are two types of logical channels in NB-IoT, the data channels and the signaling channels.
            \begin{itemize}
                \item The data channels are used to transmit the payload between the UE and the eNB.
                \item The signaling channels are used to transmit control information between the UE and the eNB, such as the messages needed to join the network, the Tracking Area Updates, etc. Among those, the Paging messages are also sent through the signaling channels.
            \end{itemize}
        \end{enumerate}
    \end{ejercicio}

    \newpage

    \begin{ejercicio}[20 points]
        In the lecture you learned about the application layer protocol MQTT.
        \begin{enumerate}[label=\alph*)]
            \item (4 points) Describe a scenario in which you would apply the retain flag of MQTT. Justify your answer.
            
            A sensor does not read analog, but digital values. Instead of sending the same value multiple times, it only sends a message when the value changes. If a subscriber subscribes to that topic after the last message was sent, it will not receive any message until the next change. To solve this issue, we can use the retain flag, which will make the broker store the last message sent (which represents the current state of the sensor) and send it to any new subscriber that subscribes to that topic.
            \item Given an MQTT configuration with four clients C1, C2, C3, C4 and a broker B:
            \begin{enumerate}[label=\roman*)]
                \item (3 points)
                \begin{itemize}
                    \item C1 publishes a message to topic \verb|basement/temperature|,
                    \item C2 publishes a message to topic \verb|groundfloor/kitchen/temperature|,
                    \item C3 publishes a message to topic \verb|secondfloor/bathroom/temperature|.
                \end{itemize}
                C4 should subscribe to which topic to receive all messages from C1, C2 and C3?


                Since the wildcard \verb|#| must be the last character in the topic, C4 should subscribe to \verb|#|, which will match all topics.
                \item (3 points)
                \begin{itemize}
                    \item C1 publishes a message to topic \verb|bathroom/ceiling/temp|,
                    \item C2 publishes a message to topic \verb|bathroom/floor/temp|,
                    \item C3 publishes a message to topic \verb|bathroom/boiler/outlet/temp|.
                \end{itemize}
                C4 should subscribe to which topic to get the environment temperature of the bathroom?
                \begin{observacion}
                    We do not want to receive any other types of measurements, e.g. humidity.
                \end{observacion}
                
                Since only the environment temperature of the bathroom is needed, the message from C3 should not be received. Therefore, C4 should subscribe to the topic \verb|bathroom/+/temp|, which will match the topics \verb|bathroom/ceiling/temp| and \verb|bathroom/floor/temp|, but not \verb|bathroom/boiler/outlet/temp|.
                \item (3 points) C1, C2, C3 subscribe to shared subscriptions \verb|\$share/exam/iot/temp|. To which topic should C4 publish, such that C1, C2, or C3 will receive the message?
                
                In this case, \verb|exam| is the group name, and \verb|iot/temp| is the topic. Therefore, C4 should publish to the topic \verb|iot/temp|, so that one of the clients C1, C2, or C3 will receive the message. Since it is a shared subscription, only one of the clients will receive the message.
            \end{enumerate}
            \item (3 points) MQTT v5 introduced multiple new features, among them ``Topic Alias''. While creating a ``Topic Alias'', why should the central broker store the topic alias mapping?
            
            There are several reasons why the central broker should store the topic alias mapping:
            \begin{itemize}
                \item The alias is connection-specific, so the broker needs to store the mapping for each client (connection).
                \item When receiving a message with a topic alias, it may match several topics (depending on the wildcard subscriptions). Therefore, the broker needs to know the original topic to match it with the subscriptions.
                \item If a publish message is sent with a topic alias and the retain flag is set, the broker needs to store the original topic to be able to send it to new future subscribers.
            \end{itemize}
            
            \item (4 points) Client C2 subscribes to the topic \verb|a/b/c| with QoS level $0$. Afterwards, client C1 publishes a message to the topic \verb|a/b/c| with QoS level $2$. Looking at the message headers: what is the QoS level between C1 and the broker? What is the QoS level between the broker and C2?
            
            The message is published with QoS level $2$, so the QoS level between C1 and the broker is $2$. However, since C2 subscribed to the topic with QoS level $0$, the broker downgrades the QoS level to $0$ when sending the message to C2. Therefore, the QoS level between the broker and C2 is $0$.
        \end{enumerate}
    \end{ejercicio}

    \newpage

    \begin{ejercicio}[14 points]
        ~
        \begin{enumerate}[label=\alph*)]
            \item\label{itm:robot}
                (6 points) A robot moves linearly with speed $v = \unitfrac[2]{m}{s}$. It starts at time $t = 0$ at location $m_0 = \unit[0]{m}$. We want it to send us its location with the ``Distance-based Reporting'' protocol from the lecture with an accuracy of $\text{acc} = \unit[10]{m}$ and max speed $v_{\max} = \unitfrac[4]{m}{s}$. Write down the next three positions after $m_0$ ($m_1, m_2, m_3$) that the robot will report.
                \begin{enumerate}
                    \item At $t=\unit[0]{s}$, the robot is at $m_0 = \unit[0]{m}$. Since no report has been sent yet, the robot will send a report with its current location, so $m_{\text{last}} = m_0 = \unit[0]{m}$.
                    
    
                    The sleep time is calculated as follows:
                    \begin{equation*}
                        T_0 = \frac{d_{\max} - \text{distance}(m_0, m_{\text{last}})}{v_{\max}} = \frac{10 - 0}{4} = \unit[2.5]{s}
                    \end{equation*}
                    Therefore, the robot will sleep for $T_0 = \unit[2.5]{s}$, and it will wake up at $t = \unit[2.5]{s}$.
                    \item At $t=\unit[2.5]{s}$, the robot is at $m_1=m(2.5)=2\cdot 2.5 = \unit[5]{m}$. 
                    
                    The sleep time is calculated as follows:
                    \begin{equation*}
                        T_1 = \frac{d_{\max} - \text{distance}(m_1, m_{\text{last}})}{v_{\max}} = \frac{10 - 5}{4} = \unit[1.25]{s}
                    \end{equation*}
                    Therefore, the robot will sleep for $T_1 = \unit[1.25]{s}$, and it will wake up at $t = \unit[3.75]{s}$.
                    \item At $t=\unit[3.75]{s}$, the robot is at $m_2=m(3.75)=2\cdot 3.75 = \unit[7.5]{m}$.
                    
                    The sleep time is calculated as follows:
                    \begin{equation*}
                        T_2 = \frac{d_{\max} - \text{distance}(m_2, m_{\text{last}})}{v_{\max}} = \frac{10 - 7.5}{4} = \unit[0.625]{s}
                    \end{equation*}
                    Therefore, the robot will sleep for $T_2 = \unit[0.625]{s}$, and it will wake up at $t = \unit[4.375]{s}$.
                    \item At $t=\unit[4.375]{s}$, the robot is at $m_3=m(4.375)=2\cdot 4.375 = \unit[8.75]{m}$.
                    
                    The sleep time is calculated as follows:
                    \begin{equation*}
                        T_3 = \frac{d_{\max} - \text{distance}(m_3, m_{\text{last}})}{v_{\max}} = \frac{10 - 8.75}{4} = \unit[0.3125]{s}
                    \end{equation*}
                    Therefore, the robot will sleep for $T_3 = \unit[0.3125]{s}$, and it will wake up at $t = \unit[4.6875]{s}$.
                \end{enumerate}

                It should be noted however that the robot will not send a report until he reaches the maximum distance of $d_{\max} = \unit[10]{m}$ from the last reported position. Before that, he will only measure his current position and calculate the sleep time.
            \item (4 points) When the robot from Task~\ref{itm:robot} keeps moving in the same direction with the same speed for a longer time, ``Distance-based Reporting'' leads to energy inefficiencies. Can we apply ``Dead Reckoning'' to make it more efficient? Justify your answer.
            
            Yes, we can apply ``Dead Reckoning'' to make it more efficient. The robot will initially report its initial position and an expected update function that states the expected position of the robot at any given time. In this case, that expected update function is:
            $$ m(t) = m_0 + v \cdot t = 0 + 2 \cdot t = 2t $$
            
            The receiver will use that expected update function to estimate the position of the robot at any given time. The robot will only send a report when the actual position of the robot deviates from the expected position by more than $d_{\max} = \unit[10]{m}$. In this case, since the actual position of the robot is always equal to the expected position, the robot will periodically check its position, but it will never send a report, which saves energy.
            \item (4 points) Describe a scenario where ``Dead Reckoning'' leads to energy inefficiency. Justify your answer.
            
            A scenario where ``Dead Reckoning'' leads to energy inefficiency is when the robot does not correctly predict its expected position. If the robot is moving in complex or unexpected ways, it will not be able to send an accurate expected update function to the receiver. Therefore, the robot will have to send a report more frequently, which leads to energy inefficiency.
        \end{enumerate}
    \end{ejercicio}

    \newpage

    \begin{ejercicio}[11 points]
        ~
        \begin{enumerate}[label=\alph*)]
            \item In the following two IoT application scenarios artificial intelligence shall be used to solve a task. Would you deploy the AI (or a specific part of it) locally on an embedded device or in the cloud? Justify your answer.
            \begin{enumerate}[label=\roman*)]
                \item (3 points) A voice-controlled virtual assistant (such as Amazon Alexa) shall wake up from a low-energy state on recognizing the user saying a specific keyword.
                
                I would deploy the AI locally on an embedded device. The main reasons for that are:
                \begin{itemize}
                    \item Deployment in the cloud would lead to a high latency, which is not acceptable for this application scenario. The user expects the virtual assistant to wake up immediately after saying the keyword.
                    \item The virtual assistant is always listening for the keyword. Given privacy concerns, it is not acceptable to send all the audio data to the cloud for processing.
                    \item Bandwidth is also an important factor. Sending all the audio data to the cloud would require a lot of bandwidth, which is not efficient.
                \end{itemize}
                \item (3 points) Imagine a video surveillance system for intruder detection that includes a smart camera. If a suspicious object is detected in the scene, it must be analysed whether it actually is an intruder.
                
                In this case, I would deploy an hybrid approach. Constantly sending all the video data to the cloud for processing would require a lot of bandwidth, which is not efficient. However, computer vision algorithms are computationally expensive and require a lot of processing power, which is not available on embedded devices. Therefore, I would deploy a simple computer vision algorithm on the embedded device to detect suspicious objects. If a suspicious object is detected, the video data of that object would be sent to the cloud for further analysis.
            \end{enumerate}
            \item (2 points) Explain the difference between pruning and quantization of an artificial neural network.
            
            They are two different techniques to optimize an artificial neural network. However, their aim is totally different.
            \begin{itemize}
                \item Prunning is a technique to reduce the number of parameters (and consequently the number of computations) of an artificial neural network. It is done by removing some of the weights of the network.
                \item Quantization is a technique to reduce the number of bits used to represent the weights of an artificial neural network (and consequently increase the speed of the computations).
            \end{itemize}
            \item (3 points) How does magnitude-based pruning decide which weights to prune? Why can that decision lead to drastic performance drops?
            
            Magnitude-based pruning decides which weights to prune based on their absolute value. The weights with the smallest absolute value (and therefore closer to zero) are pruned, since they are supposed to have the least impact on the output of the network. However, this decision can lead to drastic performance drops, since it does not take into account the importance of the weights in the context of the network. Some weights with small absolute values may be crucial for the performance of the network, and pruning them can lead to a significant drop in accuracy. In order to avoid that, it is important to fine-tune the network after pruning, so that the remaining weights can adapt to the new structure of the network.
        \end{enumerate}
    \end{ejercicio}

\end{document}